\subsection{OBS-007. A relative second-derivative criterion for all-copy injectivity}
\textbf{Status: proved. The complete polynomial argument and source have
independent whole-text reviews. A fresh standard-library reconstruction
passed all Jacobian entries, both kernel identities, all eighty second
coefficient vectors, and the final scalar rank check.}

\paragraph{The general affine-dependence criterion.}
Let $X$ be a nonempty open subset of an affine space in characteristic zero,
and let
\[
 f=(l,g):X\longrightarrow \mathbb A^u\times\mathbb A^v
\]
be a polynomial map. Suppose $l$ is dominant and $f$ has generic
differential rank $r$. Write $c=u+v-r$. At a point $x$ on the full-rank
open set, let $W_x$ be a matrix whose $c$ rows are a basis of the left
kernel of $df_x$. For a tangent vector $a$ in $\ker dl_x$, put
\[
 H_x(a)=[t^2]f(x+ta)
\]
in affine coordinates on the source. If there are $c$ such vectors
$a_1,\ldots,a_c$ for which the square matrix
\[
               \bigl(W_xH_x(a_i)\bigr)_{i=1}^c
\tag{OBS-007.1}
\]
is invertible, and the vectors and kernels are defined on an open
neighborhood by fixed Gaussian pivots, then no nonzero polynomial
\[
                 P(l,g)=A(l)+\sum_{j=1}^v B_j(l)g_j
\tag{OBS-007.2}
\]
vanishes on the image of $f$.

\begin{proof}
Assume $P\circ f=0$. Differentiating once shows that the row vector
$w_x=dP_{f(x)}$ lies in the left kernel of $df_x$. Along $x+ta$, the
coefficient of $t^2$ in the same identity is
\[
 0=w_xH_x(a)+\tfrac12
          d^2P_{f(x)}(df_xa,df_xa).
\]
For $a\in\ker dl_x$, the first component of $df_xa$ is zero.
Because $P$ is affine in $g$, its second derivative in two pure
$g$ directions is zero. Hence $w_xH_x(a)=0$ for every chosen $a$.
Writing $w_x$ in the row basis $W_x$, invertibility of (OBS-007.1)
forces $w_x=0$.

The rank and determinant conditions hold on a nonempty open subset
of $X$, because the chosen pivots and vectors depend rationally on
the source coordinates. Thus the last $v$ entries of $w_x$, namely
$B_j(l(x))$, vanish on that open set. Dominance of $l$ implies that
each polynomial $B_j$ is identically zero. The original identity
then gives $A(l(x))=0$, and dominance gives $A=0$ as well. No
claim about a single fiber being connected or a projective image
being normal is needed.
\end{proof}

\paragraph{Arithmetic certification of the criterion.}
The maps below are defined over $\mathbb Z$. A nonzero minor modulo
a prime $p$ supplies a pivot open on which the relevant linear
systems can be solved over a localization with denominators prime
to $p$. A global characteristic-zero upper bound $r$ for the rank,
together with a nonzero rank-$r$ minor modulo $p$, fixes the generic
rank exactly. All larger minors vanish as integer polynomial
identities. The left kernel is therefore locally free of rank $c$
on the chosen pivot open, and the kernel of the full-row-rank lower
map is likewise free there.

Each recorded finite-field lower-kernel vector lifts by fixing its
free coordinates as integer constants and solving for its pivot
coordinates. A recorded basis of the finite-field left kernel
lifts after an invertible constant change of a Gaussian kernel
basis. Hence the determinant in (OBS-007.1) has a rational-function
lift whose denominator is a unit at $p$ and whose residue is the
recorded nonzero determinant. It is nonzero in characteristic zero.
In particular, merely treating modular kernel vectors as constant
integer kernel vectors would be unjustified; the pivot lifting is
part of this argument.

\paragraph{The determinant map at size twenty-three.}
Take $X=(\operatorname{Mat}_{23})^3$ and
\[
 F_A(x_0,x_1,x_2,x_3)
       =\det(x_0I_{23}+x_1A_1+x_2A_2+x_3A_3).
\]
The coefficient of $x_0^{23}$ is one. Let $l$ contain all ordinary
coefficients with total last-three degree from1 through16, and
let $g$ contain those of degree17. Their numbers are
\[
 u=\binom{19}{3}-1=968,\qquad
 v=\binom{19}{2}=171,\qquad u+v=1139.
\]
The normalized determinant image has dimension at most
$2\cdot23^2+1=1059$, by CUE\_GCT G-06 and its scalar normalization.
Equivalently, on the open set of determinant pencils with invertible
constant matrix, left multiplication by that matrix's inverse gives
the displayed normalized pencil. The coefficient projection cannot
increase this dimension. The source has dimension1587.

The exact modular data use $p=101$, $A_1=\operatorname{diag}(1,\ldots,23)$,
and the two matrices generated successively from seed20261122 by
\[
 s\leftarrow1664525s+1013904223\pmod{2^{32}},
 \qquad \text{entry}=\lfloor s/2^{16}\rfloor\pmod{101}.
\]
The full Jacobian is $1139\times1587$. Its lower968 rows have
rank968; the full matrix has rank1059. Thus $c=80$.
The archives contain80 left-kernel rows,80 lower-kernel directions,
and all $80\cdot1139$ second-coefficient entries. The independent
reconstruction gives rank80 in (OBS-007.1), with nonzero determinant1
modulo101. The lower Jacobian has selected maximal minor82 modulo101.
A fresh complete check, without reusing previously verified directions,
passed in131.203 seconds.

\paragraph{Independent reconstruction.}
Write $Q_j$ for the homogeneous degree-$j$ coefficient of the
adjugate of $I+M$, with $M=x_1A_1+x_2A_2+x_3A_3$ and $Q_0=I$.
The elementary coefficient $e_j$ and adjugate obey
\[
 e_j=\frac1j\operatorname{tr}(MQ_{j-1}),\qquad
 Q_j=e_jI-MQ_{j-1}.
\tag{OBS-007.3}
\]
The Jacobian entries are the appropriate transposed coefficients
of $Q_{j-1}$. The exploratory builder differentiates (OBS-007.3)
as a full truncated series through order two in a source parameter.
The independent checker instead differentiates the adjugate only
once, and uses
\[
 [t]e_j(M+tV)=\operatorname{tr}(VQ_{j-1}),\qquad
 2[t^2]e_j(M+tV)=\operatorname{tr}(V[t]Q_{j-1}(M+tV)).
\tag{OBS-007.4}
\]
Both identities follow by differentiating the determinant. All
divisions used are by integers less than101.

The checker uses integer digit products for matrix multiplication.
With base $2^{18}$, a product of a forward-encoded row and a
reverse-encoded column has the matrix dot product as its middle
digit, because every convolution coefficient is at most
$23\cdot100^2<2^{18}$. It reconstructs the Jacobian, checks the
two kernel identities and both Jacobian ranks, recomputes all
second coefficients by (OBS-007.4), and performs exact scalar-field
elimination on the final80-square matrix. Source and input hashes
bind any resumed prefix to precisely the same computation. The
default invocation is a fresh full check:
\begin{quote}\ttfamily\small
python work/root/scripts/check\_relative\_second\_jets.py
\end{quote}
No NumPy is imported by that checker. The separate probe is an
accelerated construction, not the independent certificate.

\paragraph{The all-copy consequence.}
For every four-row $\lambda\vdash23d$ with $\lambda_2\le17$,
the certificate and criterion prove
\[
                  K_{23}(\lambda)=a_\lambda,
       \qquad q_\lambda=a_\lambda\quad\text{at }(m,n)=(17,23).
\tag{OBS-007.5}
\]
Here is the complete representation-theoretic deduction. By S-001,
a highest-weight vector uses only polars through order17.
Localization at $F_0$ and centering identify its fixed outer-degree
space injectively with highest-weight polynomials in
$C_2,\ldots,C_{17}$, as in S-007. The symmetric square
\[
 \operatorname{Sym}^2(\operatorname{Sym}^{17}\mathbb C^3)
       =\bigoplus_{i=0}^{8}S_{(34-2i,2i)}\mathbb C^3
\]
has first row at least18. Polynomial tensor products cannot decrease
the first row, and higher powers are quotients of this symmetric
square times further polynomial factors. Therefore every centered
highest-weight polynomial with first tail part at most17 is affine
in $C_{17}$. On $F_0=1$, centering changes $C_{17}$ by an invertible
linear normalization of the ordinary order17 coefficient tensor,
plus a polynomial in lower coefficients. Thus the vector has the
form (OBS-007.2).

If it vanished on the determinant locus, the criterion would make
that polynomial zero. Injectivity of the fixed-degree specialization
would make the original highest-weight vector zero. This proves
injectivity on the entire highest-weight space in every admissible
degree, hence $K_{23}=a$. S-007 already proves $q=a$ for these labels.
In particular the four degree-six candidates with multiplicities
$450,1332,2392,3577$ are all excluded by (OBS-007.5). Any remaining
four-row obstruction at $(17,23)$ must have $\lambda_2\ge18$.

\paragraph{Geometric formulation.}
The left Jacobian kernel is the conormal bundle on the smooth rank
open of the coefficient image. The tested map is the second
fundamental form restricted to tangent directions annihilating the
lower coefficient projection. Its full normal span excludes an
equivariant coherent ideal section that is affine along the last
coefficient bundle. The proof uses the explicit differential and
localization, with the group-scheme, Schur-functor and coherent-ideal
dictionary GD-1--GD-3. The arithmetic pivot opens give actual locally
free kernel modules; no general assertion that an arbitrary kernel
commutes with every base change is used. No purity or motivic input
is needed.
